% Generisano: uređuj beleske/sNNN.tex i naslove u manifestu.
\documentclass[11pt,a4paper]{article}
% Zajednički preambulum za dokument za čitanje; nije Beamer tema.
\usepackage[a4paper,margin=20mm,headheight=15pt]{geometry}
\usepackage{lmodern,fontspec}
\usepackage[serbian]{babel}
\setmainfont{Latin Modern Sans}
\setsansfont{Latin Modern Sans}
\setmonofont{Latin Modern Mono}
\usepackage{amsmath,amssymb,mathtools,graphicx,xcolor}
\usepackage{array,booktabs,tabularx,multirow,makecell,listings,fancyhdr}
\pagestyle{fancy}\fancyhf{}
\newcommand{\BeleskeTekuciID}{NEOZNACEN}
\fancyhead[L]{\small Beleške uz slajd \texttt{\expandafter\detokenize\expandafter{\BeleskeTekuciID}}}
\fancyfoot[R]{\small\thepage}
\setlength{\parindent}{0pt}\setlength{\parskip}{6pt}
\newwrite\BeleskeMapa
\immediate\openout\BeleskeMapa=\jobname.notes-map.tsv
\AddToHook{shipout/before}{%
  \immediate\write\BeleskeMapa{\BeleskeTekuciID\space\thepage}}
\AddToHook{enddocument/afterlastpage}{\immediate\closeout\BeleskeMapa}
\newcommand{\BeleskeID}[1]{%
  \def\BeleskeZadatiID{#1}%
  \ifx\BeleskeTekuciID\BeleskeZadatiID\else
    \PackageError{beleske}{Pogresan ID beleske: #1}{Proveri mapiranje slajda.}%
  \fi}
\newcommand{\BeleskaSlajda}[4]{%
  \clearpage\gdef\BeleskeTekuciID{#1}%
  {\Large\bfseries Slajd \textnormal{\texttt{\detokenize{#1}}}: #3\par}\medskip
  \includegraphics[page=#2,width=\linewidth]{\BeleskePrezentacija}\par\medskip
  \input{#4}}

\newcommand{\BeleskePrezentacija}{build/02_sinteza_kombinacionih_mreza.pdf}
\begin{document}
\BeleskaSlajda{02-s001}{1}{Sinteza kombinacionih mreža}{beleske/s001.tex}
\BeleskaSlajda{02-s002}{2}{Šta znači „trenutno“ stanje ulaza?}{beleske/s002.tex}
\BeleskaSlajda{02-s002b}{3}{Promena ulaza i kašnjenje logičkog kola}{beleske/s002b.tex}
\BeleskaSlajda{02-s003}{4}{Kašnjenje invertora pri promeni signala}{beleske/s003.tex}
\BeleskaSlajda{02-s004}{5}{Srednje kašnjenje i kratak impuls}{beleske/s004.tex}
\BeleskaSlajda{02-s005}{6}{Maksimalna specificirana kašnjenja}{beleske/s005.tex}
\BeleskaSlajda{02-s006}{7}{Kada je izlaz kombinacione mreže važeći?}{beleske/s006.tex}
\BeleskaSlajda{02-s007}{8}{Sinteza i kriterijumi minimizacije}{beleske/s007.tex}
\BeleskaSlajda{02-s008}{9}{Integrisana kola, čip i DIP kućište}{beleske/s008.tex}
\BeleskaSlajda{02-s009}{10}{Literal, proizvod i zbir}{beleske/s009.tex}
\BeleskaSlajda{02-s009b}{11}{Indeksi članova i oblici logičke funkcije}{beleske/s009b.tex}
\BeleskaSlajda{02-s010}{12}{Koraci sinteze kombinacione mreže}{beleske/s010.tex}
\BeleskaSlajda{02-s011}{13}{Od zahteva do funkcionalne tabele}{beleske/s011.tex}
\BeleskaSlajda{02-s012}{14}{Skraćena funkcionalna tabela}{beleske/s012.tex}
\BeleskaSlajda{02-s013}{15}{Zbir proizvoda iz redova sa F=1}{beleske/s013.tex}
\BeleskaSlajda{02-s014}{16}{Kanonička disjunktivna forma}{beleske/s014.tex}
\BeleskaSlajda{02-s015}{17}{Uslov za red sa logičkom nulom}{beleske/s015.tex}
\BeleskaSlajda{02-s015b}{18}{Povezivanje uslova u proizvod zbirova}{beleske/s015b.tex}
\BeleskaSlajda{02-s016}{19}{Drugi put: komplementarna funkcija}{beleske/s016.tex}
\BeleskaSlajda{02-s017}{20}{Kanonička konjunktivna forma}{beleske/s017.tex}
\BeleskaSlajda{02-s018}{21}{Indeksi potpunih proizvoda i zbirova}{beleske/s018.tex}
\BeleskaSlajda{02-s019}{22}{Šema funkcije F u obliku zbira proizvoda}{beleske/s019.tex}
\BeleskaSlajda{02-s020}{23}{Invertori i opterećenje izvora signala}{beleske/s020.tex}
\BeleskaSlajda{02-s021}{24}{Rasterećenje ulaza}{beleske/s021.tex}
\BeleskaSlajda{02-s022}{25}{Šema funkcije F u obliku proizvoda zbirova}{beleske/s022.tex}
\BeleskaSlajda{02-s023}{26}{Direktna realizacija kombinacione mreže}{beleske/s023.tex}
\BeleskaSlajda{02-s023b}{27}{Poređenje resursa dve direktne realizacije}{beleske/s023b.tex}
\BeleskaSlajda{02-s024}{28}{Dualnost i De Morganova pravila}{beleske/s024.tex}
\BeleskaSlajda{02-s024b}{29}{Primer dualnosti: transformacija u NI mrežu}{beleske/s024b.tex}
\BeleskaSlajda{02-s025}{30}{NI kola su funkcionalno potpuna}{beleske/s025.tex}
\BeleskaSlajda{02-s026}{31}{NILI kola su funkcionalno potpuna}{beleske/s026.tex}
\BeleskaSlajda{02-s027}{32}{Šta se menja NI/NILI transformacijom?}{beleske/s027.tex}
\BeleskaSlajda{02-s028}{33}{Realizacija funkcije pomoću dvoulaznih kola}{beleske/s028.tex}
\BeleskaSlajda{02-s029}{34}{Uravnotežena realizacija petoulazne I funkcije}{beleske/s029.tex}
\BeleskaSlajda{02-s030}{35}{Oprez: NI kaskade nisu obične I kaskade}{beleske/s030.tex}
\BeleskaSlajda{02-s031}{36}{Troulazno NI kolo iz dvoulaznih}{beleske/s031.tex}
\BeleskaSlajda{02-s032}{37}{Minimizacija logičkih funkcija}{beleske/s032.tex}
\BeleskaSlajda{02-s032b}{38}{Minimizacija logičkih funkcija — primer}{beleske/s032b.tex}
\BeleskaSlajda{02-s033}{39}{Karnoova karta za ručnu minimizaciju}{beleske/s033.tex}
\BeleskaSlajda{02-s034}{40}{Karnoove karte jedne i dve promenljive}{beleske/s034.tex}
\BeleskaSlajda{02-s035}{41}{Karnoove karte za tri promenljive}{beleske/s035.tex}
\BeleskaSlajda{02-s036}{42}{Karnoova karta za četiri promenljive}{beleske/s036.tex}
\BeleskaSlajda{02-s037}{43}{Karnoova karta za pet promenljivih}{beleske/s037.tex}
\BeleskaSlajda{02-s038}{44}{Karnoova karta za šest promenljivih}{beleske/s038.tex}
\BeleskaSlajda{02-s039}{45}{Površine reda k u Karnoovoj karti}{beleske/s039.tex}
\BeleskaSlajda{02-s040}{46}{Susedstvo preko ivica i između slojeva}{beleske/s040.tex}
\BeleskaSlajda{02-s041}{47}{Algoritam minimizacije}{beleske/s041.tex}
\BeleskaSlajda{02-s042}{48}{Neodređena stanja i stalni literali}{beleske/s042.tex}
\BeleskaSlajda{02-s043}{49}{Primer: minimalni zbir proizvoda}{beleske/s043.tex}
\BeleskaSlajda{02-s044}{50}{Primer: minimalni proizvod zbirova}{beleske/s044.tex}
\BeleskaSlajda{02-s045}{51}{Poređenje po broju ulaza kola}{beleske/s045.tex}
\BeleskaSlajda{02-s046}{52}{Ista funkcija: drugačiji broj čipova}{beleske/s046.tex}
\BeleskaSlajda{02-s047}{53}{Grupisanje proizvoda i broj čipova}{beleske/s047.tex}
\BeleskaSlajda{02-s048}{54}{NI realizacija: tri različite vrste čipova}{beleske/s048.tex}
\BeleskaSlajda{02-s048b}{55}{Dvoulazna NI kola: tri ili dva čipa?}{beleske/s048b.tex}
\BeleskaSlajda{02-s049}{56}{NILI realizacija: tri različite vrste čipova}{beleske/s049.tex}
\BeleskaSlajda{02-s049b}{57}{Dvoulazna NILI kola i složena CMOS realizacija}{beleske/s049b.tex}
\BeleskaSlajda{02-s050}{58}{Lažne nule i jedinice — glič}{beleske/s050.tex}
\BeleskaSlajda{02-s051}{59}{Statički hazard: izlaz treba da ostane 1}{beleske/s051.tex}
\BeleskaSlajda{02-s052}{60}{Minimizacija i prelaz između površina}{beleske/s052.tex}
\BeleskaSlajda{02-s053}{61}{Povezivanje površina: uklanjanje lažne nule}{beleske/s053.tex}
\BeleskaSlajda{02-s054}{62}{Hazardi i usklađivanje kašnjenja}{beleske/s054.tex}
\end{document}
